% NOVOS COMANDOS PARA TÓPICOS E CORES
\newcommand{\espaco}{\hspace{0.2cm}}
\newcommand{\espacolongo}{\hspace{1cm}}
\newcommand{\derivada}{\frac{d}{dx}}
\newcommand{\seta}{$\rightarrow\,\,\,$}
\newcommand{\Seta}{$\Rightarrow\,\,\,$}
\newcommand{\obs}{\color{red} Obs: \color{black}}
\newcommand{\linha}{\noindent\rule{20cm}{0.4pt}}
\newcommand{\ponto}{\textbullet \hspace{0.1cm}}
\newcommand{\cor}[2]{\color{#1}#2\color{black}}
\newcommand{\cors}[2]{\color{#1}#2 \color{black}}
\newcommand{\quadrado}{\scalebox{0.8}{$\,\, \blacksquare \,\,$ }}

% COMANDOS PARA CAIXAS COLORIDAS E TEXTOS ESTÉTICOS
\newcommand{\blueText}[1]{
    \color{DeepSkyBlue2} #1 \color{black}
}
\newcommand{\ex}[1]{
    \colorbox{DeepSkyBlue2}{\color{white}\textbf{EXEMPLO #1}\color{black}}
}
\newcommand{\legenda}[1]{
    \colorbox{DeepSkyBlue2}{\color{white}\textbf{#1}\color{black}}
}
\newcommand{\legendaMatematica}[2]{
    \phantom{#2} \hfill $#1$ \hfill #2
}
\newcommand{\centermat}[1]{
    \begin{center} $#1$ \end{center}
}
\newcommand{\cvsection}[2][]{
  \bigskip \bigskip{\color{black} \LARGE \bfseries \MakeUppercase{#2}}\\[-1ex]
  \color{orange}\rule{\linewidth}{2pt} \\ \color{black}
}
\newcommand{\cvtag}[1]{
  \tikz[baseline]\node[anchor=base,draw=black!80!white,rounded corners,inner xsep=1ex,inner ysep=0.75ex,text height=1.5ex,text depth=.25ex]{\color{gray}#1\color{black}};
  \vspace{0.25ex}
}
\newcommand{\cvlang}[2]{%
    \textcolor{black}{\textbf{#1:}}
    \textcolor{blue!70!black}{\textbf{#2}}
}
\newcommand{\divider}{
    \textcolor{black!80!white!30}{\hdashrule{\linewidth}{0.6pt}{0.5ex}}\smallskip
}
\newcommand{\topico}[1]{
    \large\color{red}#1\normalsize \color{black} \\
}
\newcommand{\titulo}[1]{
    \begin{center} \Large \color{blue} #1 \normalsize \color{black} \end{center}
}
\newcommand{\fillPage}[1]{
    \color{white} a \color{black}
    \vfill  
    
    #1
    
    \vfill
    \color{white} a \color{black}
    \newpage
}
\newcommand{\modeloTeorema}[2]{
    \begin{center}
    \begin{minipage}{.8\linewidth}
        \tcbset{
            enhanced,            % ESTILO DO QUADRO
            left=5mm,            % ESPAÇAMENTO ESQUERDO
            right=5mm,           % ESPAÇAMENTO DIREITO
            top=5mm,             % ESPAÇAMENTO SUPERIOR TEXTO
            bottom=5mm,          % ESPAÇAMENTO INFERIOR TEXTO
            boxrule=1pt,         % GROSSURA DAS LINHAS DE BORDA
            fonttitle=\bfseries, % ESTILO DA FONTE DO TÍTULO
            toptitle=2mm,   % ESPAÇAMENTO SUPERIOR TITULO
            bottomtitle=2mm % ESPAÇAMENTO INFERIOR TITULO
        }
        \begin{tcolorbox}[
            colback=blue!25!white,     % COR DO FUNDO DO TEXTO
            colframe=orange!70!yellow, % COR DO FUNDO DO TÍTULO
            title=\centering \color{black} #1
            ]
            #2
        \end{tcolorbox}
    \end{minipage}
    \end{center}
}
\newcommand{\titleBox}[4]{
    \begin{center}
    \begin{minipage}{.#1\linewidth}
        \tcbset{
            enhanced, left=4mm, right=4mm,top=4mm,
            bottom=4mm,boxrule=2pt
        }
        \begin{tcolorbox}[colback=#2]
            \color{#3} \Large \bfseries \centering
            #4
        \end{tcolorbox}
    \end{minipage}
    \end{center}
}
\newcommand{\bigpar}[1]{
    \bigg( #1\bigg )
}
\newcommand{\bordaImagem}[2]{
    \efboxsetup{linecolor=black,linewidth=#1, margin=0pt}
    \efbox{#2}
}

\newcommand{\referencia}[5]{
    \begin{flushleft} \noindent \justify \small
        \hypertarget{#1}{}
        #1. #2, #3. Disponível em: #4. Acesso em: #5.
    \end{flushleft}\vspace{-1cm}
}

\newcommand{\fcol}[2]{%
  \multicolumn{1}{|m{\dimexpr #1\linewidth - 2\tabcolsep\relax}|}{\footnotesize\centering #2}%
}

\newcommand{\col}[2]{%
  \multicolumn{1}{m{\dimexpr #1\linewidth - 2\tabcolsep\relax}|}{\footnotesize\centering #2}%
}
\providecommand{\fonte}[1]{\par\vspace{4pt}{\centering\small\textbf{Fonte:} #1\par}}
\providecommand{\reservado}[2]{\fbox{\parbox[c][#1][c]{0.8\textwidth}{\centering\textcolor{red}{(#2)}}}}